% =========================================================================
%  COPY THIS FILE to start a new note.
%
%  Switch the look with the class option:
%      [dark]   Tokyo Night, matches Obsidian on screen
%      [light]  white page, same accent colours, for printing / turning in
% =========================================================================
\documentclass[dark]{hmcnote}
\usepackage{preamble}

\course{Phys 116 -- Quantum Mechanics}
\begin{document}

\lecture{1}{Spin Angular Momentum}{2026-01-21}
QM is not just a way of asking where particles are --- it is an operating
system: polarization of light, solid state phonons, particle physics, string
theory. Anywhere you are quantizing.

\section{Stern--Gerlach}

Silver is $\mathrm{Ag} = [\mathrm{Kr}]4d^{10}5s^{1}$: one unpaired electron.
If it behaved like a spinning ball of charge we would see a smear. We see two
discrete blobs, as if each atom carried angular momentum $\pm\frac{\hbar}{2}$.

We write up as $\ket{+\vec{z}}$ or $\ket{\uparrow}$, down as $\ket{-\hat{z}}$
or $\ket{\downarrow}$. Then
%
\begin{align}
  \braket{ +z | +z } &= 1, &
  \braket{ -z | +z } &= 0, &
  \braket{ +z | +x } &= \sqrt{\tfrac{1}{2}}\, e^{i\delta}.
\end{align}

\begin{note}[Why complex?]
  The inner product $\braket{ \psi | \phi } \in \C$, unlike the familiar
  $\vec{A}\cdot\vec{B} \in \R$. That extra phase freedom is exactly what
  produces interference.
\end{note}

\subsection{Operators}

A measurement with $100\%$ probability of a given result gets its own ket.
The projection $\ketbra{+z}{+z}$ picks that component out of any state:
%
\[
  \ket{\psi} = c_{+}\ket{+\hat{z}} + c_{-}\ket{-z},
  \qquad
  c_{+} = \braket{ +\vec{z} | \psi }.
\]

\begin{definition}[Expectation value]
  For an observable $\hat{A}$ in state $\ket{\psi}$,
  $\expval{\hat{A}} = \braket{ \psi | \hat{A} | \psi }$.
\end{definition}

\begin{theorem}[Born rule]
  $\abs{\braket{ \psi | \phi }}^{2}$ is the probability of finding state
  $\psi$ given state $\phi$.
\end{theorem}

\begin{proof}
  Follows from normalization $\braket{ \psi | \psi } = 1$ together with
  completeness of $\{\ket{+z}, \ket{-z}\}$.
\end{proof}

\section{Comparison table}

\begin{center}
\begin{tabular}{@{}ll@{}}
  \toprule
  \hmctablehead{Familiar 3D vectors} & \hmctablehead{Quantum state vectors} \\
  \midrule
  $\vec{i}\cdot\vec{i} = 1$        & $\braket{ +z | +z } = 1$ \\
  $\vec{i}\cdot\vec{j} = 0$        & $\braket{ -z | +z } = 0$ \\
  $\vec{A}\cdot\vec{B} \in \R$     & $\braket{ \psi | \phi } \in \C$ \\
  \bottomrule
\end{tabular}
\end{center}

\section{Your macros still work}

$\bb{R}^3$, $\bvec{F} = m\bvec{a}$, $\pardx{\psi}{t}$, $\ddx{y}{x}$,
$\norm{\bvec{v}}$, $\abs{z}$, $\curl \bvec{E}$, $\div \bvec{B} = 0$,
$\grad \phi$, $\vspan\{v_1, v_2\}$.

\begin{warning}[Watch out]
  \verb|\div| had to be redefined --- LaTeX already uses it for $\times$-style
  division. See \texttt{tex/preamble.sty}.
\end{warning}

% Images: drop them in an images/ folder next to this file, then
%   \embed{Pasted image 20260121111713.png}
% \embed[0.7\linewidth]{Pasted image 20260121111713.png}

\end{document}
